%----------------------------------------------------------------------------------------
%    PACKAGES AND THEMES
%----------------------------------------------------------------------------------------

\documentclass[aspectratio=169,xcolor=dvipsnames]{beamer}
\usetheme{SimpleDarkBlue}

\usepackage{hyperref}
\usepackage{graphicx}
\usepackage{booktabs}
\usepackage[russianb]{babel}
\usepackage{amssymb}
\usepackage{mathtools}
\usepackage{amsmath}
\usepackage{mathrsfs}

\newcommand{\R}{\mathbb{R}}
\newcommand{\norm}[1]{\left\lVert #1 \right\rVert}
\newcommand{\fro}[1]{\left\lVert #1 \right\rVert_{\mathrm F}}
\newcommand{\rank}{\operatorname{rank}}

%----------------------------------------------------------------------------------------
%    TITLE PAGE
%----------------------------------------------------------------------------------------

\title{Геометрические методы на многообразии матриц фиксированного ранга для matrix completion}
\subtitle{Второй отчет о прогрессе}
\author{Аляутдинов Карим \& Исмаилов Магомед}
\institute{НИУ ВШЭ}
\date{11 мая 2026}

%----------------------------------------------------------------------------------------

\begin{document}

\begin{frame}
    \titlepage
\end{frame}

\begin{frame}{План выступления}
    \small
    \begin{enumerate}
        \item Что уже выполнено по проекту и в какой последовательности.
        \item Этап 1: результаты и ограничения vanilla RGD.
        \item Формулировка проблемы и мотивация перехода к compact RGD.
        \item Этап 2: реализация compact RGD и первые сравнительные результаты.
        \item Текущие ограничения и план до 01.06.2026.
    \end{enumerate}
\end{frame}

%------------------------------------------------
\section{План проекта и статус}
%------------------------------------------------

\begin{frame}{Задачи проекта (01.02.2026 -- 01.06.2026)}
    \footnotesize
    \begin{enumerate}
        \item \textcolor{green!70!black}{Формализация постановки и метрик (01.02 -- 15.02) -- выполнено}
        \item \textcolor{green!70!black}{Реализация baseline-методов ALS, Soft-Impute, RGD (16.02 -- 10.03) -- выполнено}
        \item \textcolor{green!70!black}{Реализация compact RGD + L2 (11.03 -- 05.04) -- выполнено}
        \item \textcolor{green!70!black}{Разработка API и ноутбука для массовых экспериментов (06.04 -- 30.04) -- выполнено}
        \item \textcolor{orange}{Серии синтетических экспериментов: size/rank/init/noise (01.05 -- 01.06) -- в процессе}
        \item Тест на макроэкономическом датасете (02.06 -- .05)
        \item Финализация текста, графиков и выводов (29.05 -- 01.06)
    \end{enumerate}
    \vspace{0.3cm}
    \textbf{Цель этапа:} показать, что реализован воспроизводимый пайплайн и получены первые сравнительные результаты для RGD и compact RGD.
\end{frame}

%------------------------------------------------
\section{Сделанная работа}
%------------------------------------------------

\begin{frame}{Что сделано в коде}
    \begin{itemize}
        \item Разработана архитектура быстрого проведения синтетических экспериментов
              \begin{itemize}
                  \item разделение экспериментат на этапы, за каждый из которых отвечает отдельная конфигурация
                  \item собрка отдельных конфигураций в сценарий
                  \item быстрое создание и добавление исследуемых методов
                  \item построение последовательности сценариев для исследования поведения методов при изменении одного (нескольких) параметра (ов)
              \end{itemize}
        \item Реализован модуль \texttt{matrix\_completion\_methods.py}:
              \begin{itemize}
                  \item \texttt{solve\_als}, \texttt{solve\_soft\_impute}, \texttt{solve\_riemannian\_gradient\_descent}, \texttt{solve\_riemannian\_gradient\_descent\_compact}.
              \end{itemize}
        \item Реализован \texttt{synthetic\_api.py}:
              \begin{itemize}
                  \item генерация матриц, масок, разбиений, шума;
                  \item запуск методов, агрегация и экспорт результатов.
              \end{itemize}
        \item Собран управляющий ноутбук \texttt{synthetic\_control\_panel.ipynb}:
              \begin{itemize}
                  \item one-factor sweeps, прогресс-бар \texttt{tqdm}, графики, CSV/LaTeX экспорт.
              \end{itemize}
    \end{itemize}
    % \vspace{0.2cm}
    % \textbf{Факт:} на текущем этапе вся экспериментальная часть воспроизводима через seed и единый протокол.
\end{frame}

\begin{frame}{Экспериментальный протокол}
    \small
    \begin{itemize}
        \item Синтетические матрицы: контролируем \(m,n,\rank(X_*),\) спектр сингулярных чисел, распределение элементов, когерентность.
        \item Паттерны пропусков: random / block / clustered / row-focus / column-focus.
        \item Разбиение наблюдений: fit / validation / test.
        \item Метрики: \texttt{runtime\_sec}, \texttt{iterations}, \texttt{test\_rmse}, \texttt{test\_mae}, \texttt{relative\_fro\_error}.
        \item Усреднение: несколько seed на каждую точку свипа.
    \end{itemize}
    % \vspace{0.2cm}
    % \textbf{Принцип честного сравнения:} одинаковые сценарии и маски для всех методов в каждой точке.
\end{frame}

\begin{frame}{Исследование vanilla RGD}
    \small
    Реализовали риманов градиентный спуск и провели несколько эспериментов:
    \begin{itemize}
        \item Проверили скорость сходимости алгоритма в зависимости от размера квадратной матрице (ранг фиксирован)
        \item Проверили зависимость ошибки (RMSE и MAE) от размеров матрицы
        \item Изучили скорость сходимости и ошибки в зависимости от истинного ранга матрицы
        \item Сравнили с методами ALS и Soft-Impute
    \end{itemize}
\end{frame}
\begin{frame}{Исследование vanilla RGD}
    \small
    %.                        GRAPHIC                        .%
    \textbf{График 1}
    \noindent
    \begin{center}
        \includegraphics[width=.7\textwidth]{graphics/vanilla_rgd_runtime_size-sweep.png}
    \end{center}
    %.                        GRAPHIC                        .%

\end{frame}
\begin{frame}{Исследование vanilla RGD}
    \small
    %.                        GRAPHIC                        .%
    \textbf{График 2}
    \noindent
    \begin{center}
        \includegraphics[width=0.9\textwidth]{graphics/vanilla_rgd_error_size-sweep.png}
    \end{center}
    %.                        GRAPHIC                        .%

\end{frame}
\begin{frame}{Исследование vanilla RGD}
    \small
    %.                        GRAPHIC                        .%
    \textbf{График 3}
    \noindent
    \begin{center}
        \includegraphics[width=0.7\textwidth]{graphics/vanilla_rgd_runtime_rank-sweep.png}
    \end{center}
    %.                        GRAPHIC                        .%

\end{frame}
\begin{frame}{Исследование vanilla RGD}
    \small
    %.                        GRAPHIC                        .%
    \textbf{График 4}
    \noindent
    \begin{center}
        \includegraphics[width=0.9\textwidth]{graphics/vanilla_rgd_error_rank-sweep.png}
    \end{center}
    %.                        GRAPHIC                        .%

\end{frame}
\begin{frame}{Исследование vanilla RGD}
    \small
    \textbf{Проблема} \\[30pt]
    Как мы видим, градиентный спуск демонстрирует хорошую устойчивость к изменеию размера матрицы и ее истинного ранга
    в терминах ошибок RMSE и MAE, однако очевидным слабым местом является скорость сходимости: на матрицах $2048\times 2048$
    среднее вермя восстановления составляет около 150 секунд, что примерно в 10 раз медленнее чем ALS.

\end{frame}
\begin{frame}{Исследование vanilla RGD}
    \small
    \textbf{Причина} \\[20pt]
    Происходит это из-за устройства алгоритма. Шаг градиентного спуска по многообразю происходит в 3 этапа:
    \begin{itemize}
        \item Высчитывание градиентна в текущей точке и проекция на касательное пространство;
        \item Шаг по антиградиетну в касательном пространства
        \item Возвращение на многообразие (ретракция)
    \end{itemize}
    Ретракция на многообразие матриц ранга $r$ происходит через подсчет $SVD$, что серьезно вредит
    скорости алгоритма.
\end{frame}
\begin{frame}{Исследование vanilla RGD}
    \small
    \textbf{Решение} \\[20pt]
    Нашли 2 подхода к решению этой проблемы:
    \begin{itemize}
        \item Compact RGD: Переписать формулу касательного приращения в форме SVD, разбив на ортогональные факторы и ядро, после чего делать ретракцию только по ядру, что эквивалентно SVD матрицы $2r\times 2r;$
        \item ALS-подобная эвристика: Рассмотреть решение задачи без ретракции: факторизовать матрицу, вычислить 2 градиента по прямоугольным факторам и делать одновременно 2 шага, после чего сворачивать матрицу обратно и считать ошибку.
    \end{itemize}
    На данном этапе мы исследуем Compact RGD.
\end{frame}

\begin{frame}{Исследование compact RGD}
    \small
    Реализовали компактный вариант риманова градиентного спуска и провели несколько эспериментов:
    \begin{itemize}
        \item Проверили скорость сходимости алгоритма в зависимости от размера квадратной матрице (ранг фиксирован)
        \item Проверили зависимость ошибки (RMSE и MAE) от размеров матрицы
        \item Изучили скорость сходимости и ошибки в зависимости от истинного ранга матрицы
        \item Добавили регуляризацию для лучшей устойчивости
        \item Сравнили показатели с vanilla RGD
        \item Изучили устойчивость алгоритма к ошибке в предполагаемом ранге и сравнили
        \item Сравнили с методами ALS и Soft-Impute
    \end{itemize}
\end{frame}
\begin{frame}{Исследование compact RGD}
    \small
    %.                        GRAPHIC                        .%
    \textbf{График 5}
    \noindent
    \begin{center}
        \includegraphics[width=.7\textwidth]{graphics/compact_rgd_runtime_size-sweep.png}
    \end{center}
    %.                        GRAPHIC                        .%

\end{frame}
\begin{frame}{Исследование compact RGD}
    \small
    %.                        GRAPHIC                        .%
    \textbf{График 6}
    \noindent
    \begin{center}
        \includegraphics[width=.7\textwidth]{graphics/compact_rgd_runtime_rank-sweep.png}
    \end{center}
    %.                        GRAPHIC                        .%

\end{frame}
\begin{frame}{Исследование compact RGD}
    \small
    %.                        GRAPHIC                        .%
    \textbf{График 7}
    \noindent
    \begin{center}
        \includegraphics[width=.7\textwidth]{graphics/compact_rgd_als_simpute_runtime_size-sweep.png}
    \end{center}
    %.                        GRAPHIC                        .%

\end{frame}
\begin{frame}{Исследование compact RGD}
    \small
    %.                        GRAPHIC                        .%
    \textbf{График 8}
    \noindent
    \begin{center}
        \includegraphics[width=.9\textwidth]{graphics/compact_rgd_als_simpute_error_size-sweep.png}
    \end{center}
    %.                        GRAPHIC                        .%

\end{frame}
\begin{frame}{Исследование compact RGD}
    \small
    %.                        GRAPHIC                        .%
    \textbf{График 9}
    \noindent
    \begin{center}
        \includegraphics[width=.7\textwidth]{graphics/compact_rgd_als_simpute_runtime_rank-sweep.png}
    \end{center}
    %.                        GRAPHIC                        .%

\end{frame}
\begin{frame}{Исследование compact RGD}
    \small
    %.                        GRAPHIC                        .%
    \textbf{График 10}
    \noindent
    \begin{center}
        \includegraphics[width=.9\textwidth]{graphics/compact_rgd_als_simpute_error_rank-sweep.png}
    \end{center}
    %.                        GRAPHIC                        .%

\end{frame}
\begin{frame}{Исследование compact RGD}
    \small
    \textbf{Наблюдение в экспериментах:} \\[10pt]
    Compact RGD оказался намного более устойчивым к увеличению размера матрицы и ее ранга,
    местами даже опережая ALS, при этом показывая сравнительно небольшую потерю точности.
    Проанализировав собранные данные мы заметили, что алгоритм выходит на плато и практически никогда не пробивает
    критерий остановы $tol=1e-6,$ из-за чего отрабатывает все отведенные ему итерации, при увеличении которых
    (как нам кажется) возникает плохая численная обусловленность core-матрицы, из-за чего не сходится SVD. \\[15pt]
    Попробуем тогда ограничить количество итераций.

\end{frame}
\begin{frame}{Исследование compact RGD}
    \small
    Проделав это получаем следующие результаты: \\[10pt]
    %.                        GRAPHIC                        .%
    \textbf{График 11}
    \noindent
    \begin{center}
        \includegraphics[width=.7\textwidth]{graphics/shitass1.png}
    \end{center}
    %.                        GRAPHIC                        .%


\end{frame}
\begin{frame}{Исследование compact RGD}
    \small
    %.                        GRAPHIC                        .%
    \textbf{График 12}
    \noindent
    \begin{center}
        \includegraphics[width=.9\textwidth]{graphics/shitass2.png}
    \end{center}
    %.                        GRAPHIC                        .%


\end{frame}
\begin{frame}{Исследование compact RGD}
    \small
    %.                        GRAPHIC                        .%
    \textbf{График 13}
    \noindent
    \begin{center}
        \includegraphics[width=.7\textwidth]{graphics/this_shit_is_ass.png}
    \end{center}
    %.                        GRAPHIC                        .%
\end{frame}
\begin{frame}{Исследование compact RGD}
    \small
    \textbf{Мини-итог:} \\[10pt]
    Таким образом compact SVD способен сходиться быстрее ALS и Soft-Impute, при этом относительная ошибка на больших матрицах
    мала.
\end{frame}

%------------------------------------------------
\section{Ограничения и дальнейшие шаги}
%------------------------------------------------

\begin{frame}{Сложности текущего этапа}
    \small
    \begin{itemize}
        \item На больших задачах часть запусков требует более устойчивого численного режима.
        \item Время прогонов больших свипов существенно растет.
        \item Пока нет полноценного блока экспериментов с реальными данными.
    \end{itemize}
    \vspace{0.3cm}
    \textbf{Что уже сделали для устойчивости:}
    \begin{itemize}
        \item safe-run запуск по одному методу (чтобы единичный сбой не ронял всю серию);
        \item логирование не только качества, но и времени/итераций/ошибок запуска;
        \item отдельные сценарии для диагностики инициализации и mis-specified rank.
    \end{itemize}
\end{frame}

\begin{frame}{План до финальной сдачи}
    \small
    \begin{enumerate}
        \item Завершить исследование compact RGD, решить проблемму плохой сходимости.
        \item Исследовать факторизационный вариант без явной ретракции.
        \item Завершить синтетические свипы: noise level, observed fraction, coherence mode.
        \item Добавить эксперимент на макроэкономическом датасете (masked holdout-оценка).
        \item Подготовить финальные графики/таблицы для текста курсовой.
    \end{enumerate}
    \vspace{0.3cm}
    \textbf{Ожидаемый итог:} четко описанные режимы, где RGD выигрывает по времени и сохраняет приемлемое качество.
\end{frame}

\begin{frame}{Итоги второго отчета}
    \small
    \begin{itemize}
        \item Реализована и протестирована рабочая вычислительная платформа для экспериментов.
        \item Получены первые содержательные результаты для vanilla RGD и compact RGD.
        \item Сформирован конкретный план доведения проекта до финальной защиты.
    \end{itemize}
    \vspace{0.6cm}
    \centering
    \Large Спасибо за внимание!
\end{frame}

\end{document}
